\section{Conclusiones}

\begin{itemize}

    \item El control de ratio de caudal funcionó correctamente para los tres
    setpoints evaluados (0.4, 0.6 y 0.8), manteniendo la relación FT11/FT10
    incluso ante la perturbación generada al reducir la apertura de EV11
    de 40\,\% a 20\,\%.

    \item La autosintonización mediante ``Optimización inicial'' de TIA Portal
    entregó ganancias adecuadas para cada lazo, evitando la necesidad de
    ajuste manual y reduciendo el tiempo de puesta en marcha del sistema.

    \item El control en cascada demostró mejor rechazo a perturbaciones que un
    controlador de nivel simple, ya que el lazo interno FC10 compensó
    variaciones en el caudal antes de que estas afectaran el nivel del
    tanque T-10.

    \item Se verificó que el lazo secundario FC10 debe ser más rápido que el
    lazo primario LC10 para garantizar la estabilidad de la cascada, lo cual
    quedó reflejado en las ganancias obtenidas por la autosintonización.

    \item La lógica de banda muerta implementada en el reto eliminó el llenado
    gradual del tanque causado por el comando residual de FC10, demostrando
    que los controladores PID industriales requieren lógicas complementarias
    para un desempeño óptimo.

    \item El uso del Trace de TIA Portal fue fundamental para validar el
    comportamiento de ambos lazos, permitiendo verificar el seguimiento de
    setpoint y analizar los transitorios de forma clara y precisa.

\end{itemize}